\begin{ledgroup}
\footnotesize
\pstart
\noindent
Die nachfolgenden Stücke N.\;\ref{32310}, N.\;\ref{60341}, N.\;\ref{32313} und N.\;\ref{49409}
belegen die zweite Phase von Leibnizens Auseinandersetzung mit \indexN{\textls{Vanni}, Giovanni Francesco 1638\textendash1709}Vannis kontroversen Thesen über die Statik,
die in der Publikation der \textit{Demonstratio geometrica regulae apud Staticos receptae} im November 1685 mündete (Textzeuge \textit{E\textsuperscript{1}} von N.\;\ref{49409}).
Für diesen Textkomplex wird aus folgenden Gründen die Gesamtdatierung Ende September bis November 1685 vorgeschlagen.
\pend
\pstart
In seinem \cite{02009}Brief an Leibniz vom 9. (19.) September 1685 (\textit{LSB} III,~4 N.\;96) berichtete 
\indexN{\textls{Pfautz} (Pfauzius), Christoph 1645\textendash1711}Pfautz 
von \indexN{\textls{Vanni}, Giovanni Francesco 1638\textendash1709}Vannis Entrüstung über die,
von Leibniz verfasste aber anonym erschienene,
kritische Bemerkung über sein \cite{01995}\textit{Specimen}
in den \textit{Acta Eruditorum} von November 1684 (N.\;\ref{55606}).
Seinem Brief legte Pfautz
ein \cite{02010}Schreiben 
\indexN{\textls{Vanni}, Giovanni Francesco 1638\textendash1709}Vannis an 
\indexN{\textls{Magliabechi} (Magliabecchi, Magliabecki, Magliabekius, Magliabequius), Antonio 1633\textendash1714}Magliabechi bei.
Leibniz nahm Vannis Schreiben (von dem er auch eine \cite{02010}Abschrift anfertigte: \textit{LSB} III,~4 N.\;97) zum Anlass für
die Ausführungen und Beweise im ersten Text dieser Gruppe, dem Konzept N.\;\ref{32310}.
Später diente N.\;\ref{32310} als Vorlage für den zweiten Teil der \textit{Demonstratio geometrica} (N.\;\ref{49409}, \refpassage{AE_1685_501-505_3a}{AE_1685_501-505_4b}), der sich mit der Beantwortung des in Vannis \textit{Specimen} vorgeschlagenen Problems einer auf zwei schiefen Ebenen ruhenden Kugel befasst.
\pend
\pstart
Die Lösung des Problems im zweiten Teil der \textit{Demonstratio} ist zugleich eine Anwendung von Leibnizens eigenständigem, im ersten Teil (\refpassage{AE_1685_501-505_16b}{AE_1685_501-505_3a}) erbrachten Beweis des bekannten Satzes (von Stevin, Galilei und Torricelli) der Zerlegung der Kräfte auf schiefen Ebenen.
Leibniz hebt die Originalität seines Beweises hervor: Er beruht ausschließlich auf geometrischen, nicht mechanischen Grundlagen.
Dessen erster bekannter Kern findet sich in der kurzen, nicht datierten Aufzeichung N.\;\ref{60341}. 
Leibniz übernimmt den Beweis aus N.\;\ref{60341} (bis in die Einzelheiten der Zeichnung) im ausführlichen und stark bearbeiteten Konzept N.\;\ref{32313} sowie im weiteren, auf N.\;\ref{32313} beruhenden Konzept, das hier als Textzeuge \textit{L} von N.\;\ref{49409} erscheint.
In N.\;\ref{32313} hat Leibniz außerdem auch die in N.\;\ref{32310} entworfene Beantwortung von Vannis Frage integriert und somit bereits beide Kerninhalte der \textit{Demonstratio geometrica} miteinander verbunden.
\pend
\pstart
Den sicheren Terminus ante quem für die Entstehung der Handschriften dieses Komplexes in der geschilderten Reihenfolge (in erster Linie N.\;\ref{32310} und N.\;\ref{60341}, dann N.\;\ref{32313}, schließlich der Zeuge \textit{L} von N.\;\ref{49409}) bietet wieder der Briefwechsel mit \indexN{\textls{Pfautz} (Pfauzius), Christoph 1645\textendash1711}Pfautz:
Bereits wenige Tagen nach Erhalt von Vannis Brief an Magliabechi reichte Leibniz (mit einem \cite{02011}Schreiben an Pfautz von Ende September 1685, \textit{LSB} III,~4 N.\;98)
seinen Aufsatz bei den \cite{01023}\textit{Acta Eruditorum} ein (wahrscheinlich eine, heute nicht mehr erhaltene, Reinschrift des Konzepts \textit{L} in N.\;\ref{49409}).
Daher muss er sämtliche Texte innerhalb der kurzen Zeitspanne verfasst haben, die hier als "Ende September" angegeben wird.
Eine Ausnahme bildet nur die Aufzeichnung N.\;\ref{60341}, deren Datierung lediglich als wahrscheinlich einzustufen ist. Mangels weiterer Belege ihrer Entstehung (etwa im Briefwechsel) und da das Wasserzeichen des weit verbreiteten Typus "Hollandia" keine sicheren Schlüsse zulässt, wird für N.\;\ref{60341} eine Entstehung im unmittelbaren Vorfeld der Konzepte \ref{32313} und N.\;\ref{49409} (Textzeuge \textit{L}) angenommen.
\pend
\end{ledgroup}
